% =====================================================================
% Structured summary (at most 3 pages), shared by Maria's PRDL-DL and
% MLLB-DL reports; the last paragraph depends on \ifprdl.
% Required by the professor's message of 20/11/2026: scope, core
% methodology, main findings, key conclusions.
% =====================================================================
\section*{Structured summary}
\addcontentsline{toc}{section}{Structured summary}

\paragraph{Scope.}
My question was how much each basic deep-learning module of the course contributes on a real
biomedical signal. The task is the multi-label recognition of the five diagnostic superclasses (NORM, MI,
STTC, CD, HYP) in PTB-XL \cite{wagner2020ptbxl,ptbxl_physionet,pollard2026physionet}: \edaNRecordsAll{}
twelve-lead ECGs of 10\,s from \edaNPatientsAll{} patients, sampled at 100\,Hz, each with one or more
superclasses. As the course guidelines ask \cite{course_message_2011}, the project spans core
architectures (a dense baseline, a 1D CNN, LSTM and GRU, RNN + attention, self-attention and a
Transformer encoder), learning strategies (a self-supervised model we pre-trained ourselves, then reused
frozen or fine-tuned, and zero-shot CLIP) and basic generative AI (two autoregressive ECG generators),
always with models that fit a 4\,GB GPU. The work was shared with Oriana Ferraz. I was in charge of the
data pipeline, the evaluation protocol, the recurrent networks and the self-attention and Transformer
study, which this report treats in depth; Oriana's blocks (exploratory analysis, dense network, CNN,
generators) are summarised and attributed; the learning-strategy experiments and the final comparison
were joint work.

\paragraph{Core methodology.}
\begin{itemize}[leftmargin=1.2em,itemsep=1pt,topsep=2pt]
  \item \textbf{Data.} Labels are built from the diagnostic statements and reproduce the published class
        counts. The dataset's own folds keep every patient in one split (training 1--8, validation 9,
        test 10) \cite{strodthoff2021deep}. The \edaNNosuperAll{} records without any superclass are not
        used for classification, the same choice as the benchmark \cite{strodthoff2021deep}, and the
        z-score of each lead uses training statistics only.
  \item \textbf{Metrics.} Threshold-free macro AUROC as the main score, per-class AUROC, and macro F1 with
        thresholds learnt on validation.
  \item \textbf{Comparisons.} The design of each block was fixed in writing before running it. Each
        ablation changes one factor, each configuration is trained with seeds 42, 43 and 44, model
        selection never looks at the test fold, and each final model is tested once. Differences are
        judged seed by seed on paired patient-bootstrap resamples; I call a difference real only when its
        95\,\% interval excludes 0 in all three seeds.
  \item \textbf{Reproducibility.} Each run leaves a JSON log, the report's numbers are macros generated
        from those logs, and the literature is cited only through quotes verified in the PDFs.
\end{itemize}

\paragraph{Main findings.} Table~\ref{tab:summary} gathers the final models (test macro AUROC, mean
$\pm$ SD over three seeds).
\begin{table}[htbp]
\centering
\caption[Summary of the final models (test)]{Final models on the test fold: test macro AUROC (mean $\pm$
SD over seeds 42--44; one run for the trivial model and for CLIP), trainable parameters and share of
training labels used. All final models: Table~\ref{tab:comparison_final}.}
\label{tab:summary}
\small
\begin{tabular}{llrcc}
\toprule
Model & Input & Parameters & Labels & Test macro AUROC \\
\midrule
Trivial (logistic regression on lead statistics) & 12 leads & \comparisonParamsTrivial{} & 100\,\% & \comparisonTestAurocMeanTrivial{} \\
Dense network & 12 leads & \comparisonParamsTwelveDense{} & 100\,\% & \comparisonTestAurocMeanTwelveDense{} $\pm$ \comparisonTestAurocSdTwelveDense{} \\
1D CNN & 12 leads & \comparisonParamsTwelveCnn{} & 100\,\% & \comparisonTestAurocMeanTwelveCnn{} $\pm$ \comparisonTestAurocSdTwelveCnn{} \\
RNN (convolutional front end + LSTM) & 12 leads & \comparisonParamsTwelveRnn{} & 100\,\% & \comparisonTestAurocMeanTwelveRnn{} $\pm$ \comparisonTestAurocSdTwelveRnn{} \\
RNN + attention (LSTM, 1000 steps) & 12 leads & \comparisonParamsTwelveRnnattn{} & 100\,\% & \comparisonTestAurocMeanTwelveRnnattn{} $\pm$ \comparisonTestAurocSdTwelveRnnattn{} \\
Transformer (no positional encoding) & 12 leads & \comparisonParamsTwelveTransformer{} & 100\,\% & \comparisonTestAurocMeanTwelveTransformer{} $\pm$ \comparisonTestAurocSdTwelveTransformer{} \\
\midrule
1D CNN & lead II & \comparisonParamsOneleadCnn{} & 100\,\% & \comparisonTestAurocMeanOneleadCnn{} $\pm$ \comparisonTestAurocSdOneleadCnn{} \\
ECG-GPT from scratch & lead II & \comparisonParamsScratchFOneZeroZero{} & 100\,\% & \comparisonTestAurocMeanScratchFOneZeroZero{} $\pm$ \comparisonTestAurocSdScratchFOneZeroZero{} \\
ECG-GPT fine-tuned & lead II & \comparisonParamsFinetuneFOneZeroZero{} & 100\,\% & \comparisonTestAurocMeanFinetuneFOneZeroZero{} $\pm$ \comparisonTestAurocSdFinetuneFOneZeroZero{} \\
ECG-GPT fine-tuned & lead II & \comparisonParamsFinetuneFOneZero{} & 10\,\% & \comparisonTestAurocMeanFinetuneFOneZero{} $\pm$ \comparisonTestAurocSdFinetuneFOneZero{} \\
CLIP zero-shot (ECG drawn as an image) & lead II & 0 & 0\,\% & \comparisonTestAurocMeanZeroshot{} \\
\bottomrule
\end{tabular}
\end{table}
\begin{itemize}[leftmargin=1.2em,itemsep=1pt,topsep=2pt]
  \item \textbf{Architectures.} Every family that shares weights over time is far above the dense network;
        the CNN is the best, and my final RNN is below it by \rnnDiffCnnMean{} and the final Transformer by
        \transformerDiffCnnMean{}, worse in all seeds.
  \item \textbf{Recurrent networks.} Reading all hidden states (mean or attention) adds
        \rnnDiffMeanMean{} / \rnnDiffMeanAttn{} to the last state, a GRU adds \rnnDiffMeanGru{} to an LSTM,
        and a convolutional front end adds \rnnDiffMeanConv{}; with its last state the LSTM depends mostly
        on the end of the record (\rnnSensLastBase{} of the input sensitivity in the last 2\,s). Attention
        spreads its weights almost uniformly (normalised entropy \rnnAttnEntropyNormAttn{}): it works as a
        learned mean, and the two final recurrent models cannot be told apart.
  \item \textbf{Self-attention.} Built one component at a time, the largest gains come from several heads
        with $W_O$ and from the residual path with LayerNorm. No Transformer used the order of the patches
        (shuffling them changes the AUROC by at most \transformerOrderDropLayersFour{}), so removing the
        positional encoding helped (\transformerDiffMeanNoPos{}) and the final model classifies a set of
        100\,ms patches.
  \item \textbf{Learning strategies.} Fine-tuning the self-supervised ECG-GPT beats training it from
        scratch (\transferTestDiffMeanFinetuneVsScratchFOneZeroZero{} with all labels,
        \transferTestDiffMeanFinetuneVsScratchFOneZero{} with 10\,\%), and with 10\,\% of the training labels
        it beats scratch with all of them, although its early stopping still used the whole validation
        fold; its frozen features are weak, and on lead II a CNN from scratch is still better. With one lead instead of twelve every family loses between \leadsTestDiffMeanRnnattn{} and
        \leadsTestDiffMeanTransformer{}. CLIP recognises ECG images but cannot diagnose them (macro AUROC
        \zeroshotTestMacroAuroc{}; below 0.5 for MI, CD and HYP).
  \item \textbf{Generation} (Oriana's block). The generators code lead II in fewer bits than a counting
        bigram; positional encoding is essential for generation, unlike for my classifier; greedy decoding
        collapses to a flat line and sampling with temperature gives realistic beats.
\end{itemize}

\paragraph{Key conclusions.}
(i) On this dataset, weight sharing over time is the main ingredient, and local features are what the ECG
rewards most: the best recurrent model has a convolutional front end, and the CNN beats a fine-tuned
self-supervised Transformer on the same lead. (ii) Long memory and token order, the mechanisms that make
recurrence and attention powerful on language, were not exploited by our models on these labels:
attention pooling acted as a mean, and no Transformer used the order of its tokens. (iii) Self-supervised pre-training makes a model
far more label-efficient, but the best generator is not the best feature extractor. (iv) Zero-shot needs a
shared space between inputs and class descriptions that covers the task. (v) Fixing the protocol before the
runs and pairing every comparison showed which differences are real and which are not. Limits: one dataset,
five labels, three seeds, small models and budgets, no external or clinical validation; interpretation
figures are observations on a few records.

\paragraph{How this report is organised.}
\ifprdl
Section~\ref{sec:data} describes the problem and my data pipeline, Section~\ref{sec:evaluation} the
evaluation protocol, and Section~\ref{sec:architectures} the core architectures, with the recurrent,
attention and Transformer models in depth. Sections~\ref{sec:strategies} and~\ref{sec:generative} cover the
learning strategies and the generative models, Section~\ref{sec:discussion} the final comparison, the
published results and the conclusions, and Section~\ref{sec:activities} my activities.
Appendix~\ref{app:ai} documents the use of AI assistants and Appendix~\ref{app:strategies} gives my main
takeaways on learning strategies.
\else
Section~\ref{sec:ai-env} describes the AI environment, Section~\ref{sec:methodology} the experimental
method, Section~\ref{sec:mllb-data} the data pipeline in PyTorch and Section~\ref{sec:mllb-training} the
training loop and the evaluation code. Section~\ref{sec:mllb-arch} explains the PyTorch implementation of
my recurrent and attention models, Section~\ref{sec:mllb-strategies} the implementation of the learning
strategies and of the generator, Section~\ref{sec:code} the code distribution and how to reproduce every
result, and Section~\ref{sec:mllb-conclusions} the conclusions. Appendix~\ref{app:ai-code} documents the
errors of the coding assistant.
\fi
